\section{First Principle: Keep the Feedback That Gives a Boundary Shape}

A boundary is useful when it permits the right exchanges and resists the wrong ones. A cell membrane, a professional role, and a person's commitments differ greatly, yet each depends on a distinction that has to survive contact.

The chapter calls this a \emph{damping ontology}: understand a practical boundary partly through how it receives, limits, and answers disturbance. The phrase brings attention to interaction. It supplies no single physical law for membranes, identities, and institutions.

\subsection{A Boundary Has to Be Maintained}

An illustrative balance writes
\begin{equation}
 \dot B=-\frac{B}{\tau}+P.
\end{equation}
$B$ is a defined measure of a maintained capability or boundary, $\tau$ its decay timescale in this model, and $P$ the supporting input per unit time. With constant $P$ and positive $\tau$, the balance tends toward $B^\ast=P\tau$.

The equation describes a mechanism to test. Some boundaries persist passively; others need substantial continuing work. Psychological or institutional variables must be defined and measured before this model can explain their change.

Living systems offer a concrete maintenance account. Recognition, tolerance, repair, and controlled exchange preserve an organism's functioning \cite{burnet1959clonal}. An immune system's distinction is more elaborate than a simple self/nonself label, and neither all mistakes nor all changes are fatal.

Human boundaries develop through relationships, action, learning, and institutions. A disagreement can make a preference explicit. Keeping a promise can establish reliability. A usable right of refusal can keep a relation voluntary. These interactions need continuity and response.

The functional aim is selective exchange. A sealed wall and an unprotected opening are both poor models of a capable living boundary.

\subsection{Why Some Signals Need a Stake}

Zahavi's handicap account and later signaling models explain how different costs can support honest distinctions \cite{zahavi1975mate,grafen1990biological}. A simple separating case is
\begin{equation}
 c(q_{\mathrm{low}})>b>c(q_{\mathrm{high}}),
\end{equation}
where $b$ is a benefit and $c(q)$ the cost of sending a signal for type $q$. Under the model's incentive assumptions, the higher type benefits from signaling while the lower type does not.

The relation between cost, benefit, type, and possible deception does the work. A signal need not always impose a large equilibrium cost to remain informative; direct indices, aligned interests, enforcement, and differential incentives can also sustain honesty.

This precision matters when AI makes polished artifacts cheap. An article, portfolio, or résumé may reveal less about an individual's unaided production skill. It can still reveal choice, judgment, taste, or effective use of a tool if the assessment is designed around those properties.

Three kinds of evidence deserve particular attention. A traceable history can be harder to fabricate than a claim about history. Task-specific demonstration can establish a capability more directly than an attractive sample. An enforceable stake can give an assertion consequences for the party making it.

Time contributes when the property itself develops or is observed over time. Sustained reliability requires repeated opportunities to behave reliably. Skill formation may require practice and recovery. Better training can reduce unnecessary delay, so “it takes time” should identify the relevant learning or biological process.

An expensive ordeal that measures little is a poor signal. A cheap test that directly establishes the desired capability can be an improvement.

\subsection{Formation and Maintenance Are Different Bills}

A skill must be acquired and then preserved. A record must be created and then kept readable. A relationship must begin and then continue. Creation cost and maintenance cost can move independently.

Thermodynamics gives a precise balance for a physical open system:
\begin{equation}
 \frac{\mathrm d S_{\mathrm{sys}}}{\mathrm dt}
 =\dot S_{\mathrm{prod}}+\dot S_{\mathrm{flow}},
 \qquad \dot S_{\mathrm{prod}}\geq0.
\end{equation}
Entropy has units $\mathrm{J\,K^{-1}}$ and the rates $\mathrm{J\,K^{-1}\,s^{-1}}$. A steady internal entropy requires outward entropy flow to balance internal production under the specified boundary and sign convention.

Schrödinger's discussion of life and the work on dissipative structures make the physical maintenance problem vivid \cite{schrodinger1944,nicolis1977,prigogine1977nobel}. A social practice also needs a maintenance account, but its trust or meaning cannot simply be inserted as thermodynamic entropy.

A library survives through storage, cataloging, migration, interpretation, and use. A standard survives through implementation, revision, and institutions willing to resolve disputes. The important commonality is a continuing chain of support.

Artificial tools can lower both formation and maintenance costs. They can recover a damaged text, translate it, check references, and help new readers enter a field. The result depends on whether the necessary functions remain funded and performed.

Cost does not determine worth. A costly institution can preserve something harmful; a nearly effortless act can matter deeply. The normative question is what should be maintained, for whom, and at whose expense.

The first principle is consequently functional: preserve feedback, practice, records, and enforceable commitments where they sustain a valued capability or boundary. Improve the means whenever the same function can be supplied at lower cost.
